\documentclass[11pt]{article}
\usepackage[a4paper,margin=27mm]{geometry}
\usepackage{amsmath,amssymb,amsthm,mathtools,microtype,hyperref}
\usepackage[T1]{fontenc}
\usepackage{lmodern}
\hypersetup{hidelinks}
\newtheorem{theorem}{Theorem}
\newtheorem{lemma}[theorem]{Lemma}
\newtheorem{corollary}[theorem]{Corollary}
\theoremstyle{remark}
\newtheorem{remark}[theorem]{Remark}
\newcommand{\ddc}{dd^{\mathrm c}}
\newcommand{\PP}{\mathbb P}
\title{Nearly flat residuals and strict ordinary types\newline
under blowups at poles}
\author{Research manuscript}
\date{October 8, 2026}
\begin{document}
\maketitle
\begin{center}
\small Working draft. Independent mathematical review pending.
\end{center}
\begin{abstract}
We propose a strict-current criterion for ordinary types with an effective
logarithmic divisor and a residual whose ordinary type admits arbitrarily
small negative curvature. A Gauduchon trace estimate shows that its point
Lelong numbers are automatically zero. For a smooth real background, possibly nonclosed,
such a type contains a strict current exactly when the residual background
has a smooth strictly positive representative in its smooth gauge class.
The small-curvature property persists under smooth compact modifications.
The residual leaves every ambient multiplier ideal unchanged; for a normal
crossing logarithmic divisor this gives all intrinsic ideal traces on
every smooth target not contained in that divisor, even when the naive
residual restriction is identically minus infinity.
The necessity uses Demailly's established regularization theorem on the
closed current of the residual Hessian. In a declared nonclosed ruled
family with unbounded nonanalytic residual, this excludes mutual-type
strictness after every nonempty finite composition of point blowups,
including centers in the pole set. A one-sided strict witness persists,
as do the two already computed positive ordinary-type masses. These are
restricted proposed conclusions; the general S11 programme remains open.
\end{abstract}

\section{A precise ordinary-type criterion}
Throughout $\ddc=(i/(2\pi))\partial\bar\partial$. Fix a connected compact
complex manifold $M$ and a Hermitian form $h$. Every quasi-plurisubharmonic
function has its genuine canonical representative. Write $u\sim v$ for
mutual bounded order, including values at poles. A strict current in the
ordinary type $[r]$ for a smooth real $(1,1)$-form $\theta$ means a genuine
$v\sim r$ with $\theta+\ddc v\geq\epsilon h$, $\epsilon>0$.

Let $\Delta=\sum_{j=1}^N a_jD_j$ be a finite effective real divisor,
$a_j\geq0$. Smooth Hermitian metrics on its divisor bundles give a global
logarithmic potential $\ell_\Delta=\sum_j a_j\log\|s_j\|^2$ and a smooth
closed real form $\chi_\Delta$ with
\[
 \ddc\ell_\Delta=[\Delta]-\chi_\Delta.
\]
Zero coefficients contribute the identically zero function, including
on their divisors. For $\Delta=0$, use $\ell_\Delta=\chi_\Delta=0$.

\begin{lemma}\label{lem:automatic}
Let $q$ be a genuine quasi-psh function on $M$. If for every $\delta>0$
its ordinary type contains a genuine $q_\delta$ with
$\ddc q_\delta\geq-\delta h$, then $q$ has zero point Lelong numbers
everywhere.
\end{lemma}
\begin{proof}
Write $n=\dim_{\mathbb C}M$. For $n\geq2$ choose a smooth Gauduchon
metric $\gamma$, with $\ddc\gamma^{n-1}=0$, by
\cite[Theorem~2]{Gauduchon}. That theorem and its proof apply in every
complex dimension; the even-Betti hypothesis of the source's separate
Theorem~1 is not used. For $n=1$ choose any positive $\gamma$, since
$\gamma^0=1$. The current $T_\delta=\ddc q_\delta+\delta h$ is positive,
although it need not be closed. Distributional Stokes gives its exact trace
\[
 \int_M T_\delta\wedge\gamma^{n-1}
       =\delta\int_M h\wedge\gamma^{n-1}.
\]
Fix finitely many coordinate pairs $V_a\Subset U_a$ with the $V_a$
covering $M$. On each $U_a$ take a fixed smooth positive quadratic
$\rho_a$ with $\ddc\rho_a\geq h$, and let $\beta_a$ be the Euclidean
K\"ahler form. Then $p_{\delta,a}=q_\delta+\delta\rho_a$ is a genuine
psh function and its Hessian is a closed positive local current. Uniform
comparison of the smooth metrics on these charts gives
\[
 \int_{U_a}\ddc p_{\delta,a}\wedge\beta_a^{n-1}
 \leq C_a\int_M T_\delta\wedge\gamma^{n-1}+C_a\delta
 \leq C'_a\delta.
\]
Here one can first shrink the larger charts so their closures are compact
in coordinate domains; the quadratics and all comparison constants are
fixed independently of $\delta$ and its representative. Every point of
$V_a$ has a fixed-radius coordinate ball contained in $U_a$.
The local Lelong trace bound
\cite[III.5.5--5.8]{DemaillyBook} for the closed positive Hessian therefore
bounds its point Lelong number by $C''_a\delta$. Smooth additions do not
change a point Lelong number, and mutual bounded order preserves it by
the growth formula \cite[III.6.9(c)]{DemaillyBook}. The source operator is
twice ours; its fixed positive normalization factor is absorbed into the
constants and does not affect vanishing. Thus $0\leq\nu(q,x)\leq C\delta$
for every $x$ and every $\delta>0$. Letting $\delta\downarrow0$ proves
vanishing at every point. We used closedness only for the corrected local
Hessian, not for the global Hermitian form or for $T_\delta$.
\end{proof}

\begin{corollary}\label{cor:all-modifications}
Under the hypothesis of Lemma~\ref{lem:automatic}, let
$\pi:M'\to M$ be any smooth connected compact modification. The genuine
pullback $\pi^*q$ has the same small-curvature ordinary-type property on
$M'$, and zero point Lelong numbers everywhere on $M'$.
\end{corollary}
\begin{proof}
Fix $h'>0$ on $M'$. The smooth semipositive form $\pi^*h$ is bounded
above by $C h'$ for a fixed $C>0$. Genuine quasi-psh pullbacks are
canonical and locally integrable: after a local smooth correction they
are holomorphic pullbacks of psh functions, not identically minus infinity
on an open set since a modification is generically an isomorphism.
Pointwise ordinary inequalities pull back, including their pole values.
Thus $\pi^*q_\delta\sim\pi^*q$ and
$\ddc\pi^*q_\delta\geq-C\delta h'$. Apply
Lemma~\ref{lem:automatic} on $M'$ and let $\delta$ tend to zero.
This proves the assertion for this small-curvature class; it assumes no
general preservation theorem for arbitrary zero-Lelong potentials.
\end{proof}

\section{Ambient ideals and intrinsic traces}
We use precisely $\mathcal I(tu)=\{g:|g|^2e^{-tu}\text{ is locally
integrable}\}$, without an extra factor two.

\begin{lemma}\label{lem:ideal-neutral}
If a genuine quasi-psh function $q$ has zero point Lelong numbers, then
$e^{-s q}$ is locally integrable for every $s>0$. For every genuine
quasi-psh function $v$ and every $t>0$,
\[
 \mathcal I\bigl(t(v+q)\bigr)=\mathcal I(tv).
\]
\end{lemma}
\begin{proof}
Here the needed classical integrability fact follows directly from the
Bergman approximation \cite[Theorem~1.2]{DemaillyBergman}. Locally add
a smooth correction to obtain a psh $\varphi$ with zero point Lelong
numbers on a bounded pseudoconvex coordinate domain. For any $m>0$,
the weighted Bergman potential $\varphi_m$ has
$0\leq\nu(\varphi_m,x)\leq\nu(\varphi,x)=0$.
If every weighted square-integrable holomorphic function vanished at $x$,
the mean-value and Cauchy estimates would give
$\sup_{\|g\|\leq1}|g(z)|\leq C|z-x|$ near $x$.
Then the kernel formula would imply
$\varphi_m(z)\leq m^{-1}\log|z-x|+C$ and
$\nu(\varphi_m,x)\geq m^{-1}$, a contradiction.
An actual weighted holomorphic function is therefore nonzero at $x$.
On a smaller neighborhood it is bounded away from zero, and its defining
weighted norm proves $e^{-2m\varphi}\in L^1$. Smooth corrections do not
change local integrability. Taking $2m=s$ proves the first assertion.

The local upper bound on $q$ gives
$\mathcal I(t(v+q))\subset\mathcal I(tv)$.
For the converse let $g\in\mathcal I(tv)$ at a point. Smooth local
corrections make $v$ and $q$ psh without changing any ideal. Strong
openness \cite[Corollary~2.3]{DemaillyBergman} supplies $p>1$ with
$|g|^2e^{-ptv}\in L^1$ on a smaller neighborhood.
Set $p'=p/(p-1)$. The first assertion and local boundedness of $g$ give
$|g|^2e^{-p'tq}\in L^1$. H\"older applied to the exact factorization
\[
 |g|^2e^{-t(v+q)}
   =\bigl(|g|^2e^{-ptv}\bigr)^{1/p}
      \bigl(|g|^2e^{-p'tq}\bigr)^{1/p'}
\]
proves the required inclusion. These are applications of established
Bergman approximation and strong openness, with their cited foundations;
no new integrability or openness theorem is claimed.
\end{proof}

\begin{theorem}\label{thm:intrinsic-trace}
Suppose $q$ has the small-curvature property of Lemma~\ref{lem:automatic}.
For any finite effective divisor $\Delta$, every $t>0$ satisfies
\[
 \mathcal I\bigl(t(\ell_\Delta+q)\bigr)
       =\mathcal I(t\ell_\Delta).
\]
Assume in addition that $\Delta=\sum_j a_jD_j$ has simple normal crossings.
Let $Y\subset M$ be a smooth connected compact submanifold contained in
none of the $D_j$ with $a_j>0$. For every smooth connected compact
modification $\rho:Z\to Y$ and every prime divisor $E\subset Z$, the
original intrinsic ideal trace is finite and equals
\[
 \tau_{Y,\ell_\Delta+q}(E)
 =\lim_{t\to\infty}\frac1t
   \operatorname{ord}_E\bigl(\mathcal I(t(\ell_\Delta+q))\mathcal O_Y\bigr)
 =\sum_{j:a_j>0} a_j\operatorname{ord}_E\rho^*(D_j|_Y).
\]
The assertion uses ambient ideals and imposes no nontriviality on the
naive restriction $q|_Y$.
\end{theorem}
\begin{proof}
Lemma~\ref{lem:automatic} and Lemma~\ref{lem:ideal-neutral} give the
first equality of sheaves, for arbitrary $\Delta$.
For the trace formula, normal crossing coordinates and one-variable
polar integrability give
\[
 \mathcal I(t\ell_\Delta)
       =\mathcal O_M\left(-\sum_j\lfloor t a_j\rfloor D_j\right).
\]
Indeed a holomorphic germ must vanish along each coordinate divisor to
order at least $\lfloor t a_j\rfloor$; after division by these powers,
all remaining exponents are strictly greater than $-2$ and the product
weight is integrable. At integer $t a_j$ the next smaller vanishing order
has exponent $-2$ and diverges. Smooth metric factors do not affect this.
Restriction to $Y$ is the principal ideal locally generated by
$\prod_{j:a_j>0}(g_j|_Y)^{\lfloor t a_j\rfloor}$. Zero coefficients
are omitted before restricting equations or taking orders. Every nonzero restricted
equation pulls back to a nonzero holomorphic equation on $Z$.
Its order $N_{j,E}$ along $E$ is a finite nonnegative integer.
Thus the exact order is $\sum_j\lfloor t a_j\rfloor N_{j,E}$.
Division by $t$ and convergence of the finitely many coefficients prove
the formula. No assertion that lower order constraints automatically
give an exact envelope trace is used.
\end{proof}

\section{The strict-current criterion continued}
\begin{theorem}\label{thm:criterion}
Let $q$ be a genuine quasi-psh function. Suppose that for every
$\delta>0$ there is a genuine $q_\delta\sim q$ with
$\ddc q_\delta\geq-\delta h$.
Put $B=\theta-\chi_\Delta$. Then the ordinary type
$[\ell_\Delta+q]$ contains a strict $\theta$-current if and only if
there is a smooth real $f$ with $B+\ddc f>0$.
\end{theorem}

Lemma~\ref{lem:automatic} supplies the zero-Lelong property from the
specified ordinary-type hypothesis. This does not assert that a smooth
approximation stays in $[q]$. The small-curvature representatives are
used separately for sufficiency. The regularization mechanism for necessity is established
background \cite{Demailly92}, rather than a new smoothing theorem.

\begin{lemma}\label{lem:smooth}
If $q'$ has zero point Lelong numbers and
$B+\ddc q'\geq\epsilon h$, there is a smooth real $f$ such that
$B+\ddc f\geq\epsilon h/2$.
\end{lemma}
\begin{proof}
Apply \cite[Main Theorem~1.1]{Demailly92} to the closed almost positive
current $T=\ddc q'$, its global smooth representative $0$, and the
continuous real lower bound $\gamma=\epsilon h-B$. Closedness of
$\gamma$ is not a hypothesis. The source writes
$T=(i/\pi)\partial\bar\partial\psi$, so take $\psi=q'/2$.
For any fixed attenuation parameter $s>0$, it produces global potentials
$\psi_{s,k}$ smooth off $E_s(T)$ with
\[
 \frac{i}{\pi}\partial\bar\partial\psi_{s,k}
 \geq\gamma-\min\{\lambda_k,s\}u-\epsilon_k h.
\]
Here $\lambda_k$ are continuous, decrease pointwise to the Lelong function,
and $\epsilon_k\downarrow0$. A fixed smooth tangent metric has uniformly
bounded curvature on compact $M$, so $u=Ch$ can satisfy the stronger
Griffiths curvature hypothesis used in the source's Section~2 proof.
All Lelong level sets $E_s(T)$ are empty. Dini's theorem therefore gives
$\lambda_k\to0$ uniformly. For a sufficiently large $k$, the curvature
loss is at most $\epsilon h/2$. With $f=2\psi_{s,k}$ this proves the claim.
Continuous intermediate approximants, rather than the original unbounded
$q'$, are the objects to which the source applies Richberg smoothing.
\end{proof}

\begin{proof}[Proof of Theorem~\ref{thm:criterion}]
First suppose that $v\sim\ell_\Delta+q$ is strict. Off the support
$A$ of $\Delta$, set $q'=v-\ell_\Delta$. In a local trivialization write
$\ell_\Delta=\sum_j a_j\log|g_j|^2-\sigma$, where
$\ddc\sigma=\chi_\Delta$. Choose a smooth local quadratic $\rho$
with $\ddc\rho\geq\theta$. Then
\[
 v+\rho-\sum_j a_j\log|g_j|^2=q'+\rho-\sigma
\]
is psh off $A$ and locally bounded above, since $q'\leq q+C$ and
$q$ is locally bounded above. Canonical pluripolar removability
\cite[I.5.23--5.24]{DemaillyBook} extends it across $A$.
The sharper smooth lower bound
$\ddc\rho-\theta+\epsilon h$ is retained: off $A$ it holds, while
on $A$ the coefficient matrix of the extended positive Hessian has positive
measure restriction and the smooth correction has zero restriction.
Thus the resulting global canonical residual satisfies
\[
 B+\ddc q'\geq\epsilon h.
\]
The mutual inequalities $q'\sim q$ extend canonically, by adding common
local smooth corrections and using uniqueness of psh representatives.
Lemma~\ref{lem:automatic} and bounded difference give zero point Lelong
numbers for $q'$. Lemma~\ref{lem:smooth} now gives the required smooth $f$.

Conversely, compactness gives $B+\ddc f\geq\epsilon h$ for some
$\epsilon>0$. Choose $q_\delta$ with $\delta<\epsilon/2$ and set
$v=\ell_\Delta+q_\delta+f$. The sum is a genuine quasi-psh function
and $v\sim\ell_\Delta+q$. Poincar\'e--Lelong gives
\[
 \theta+\ddc v=[\Delta]+B+\ddc f+\ddc q_\delta
                \geq\epsilon h/2.
\]
This proves sufficiency in the specified ordinary type.
\end{proof}

\section{An unbounded residual on a nonclosed ruled surface}
Let $Y=\PP^1\times\PP^1$ with probability Fubini--Study forms $w_1,w_2$.
Take $G_p\leq0$ with $\ddc G_p=[p]-w$ and $\max G_p=0$ on the first
sphere, and write
\[
 D=\{p\}\times\PP^1,\quad U=1-(1-G_p)^\alpha,\quad
 r=bG_p+U,\quad \theta=cw_1+aw_2,
 \qquad 0<\alpha<1,\quad 0\leq b<c.
\]
Here $a>0$ is smooth on the first sphere; nonconstant $a$ makes $\theta$
nonclosed. The declared residual background is
$\beta=(c-b)w_1+aw_2>0$. The already checked family \cite{Ruled}
has the two full ordinary-type masses
\[
 R_-(\theta,[r])=2(c-b)\min a,
 \qquad R_+(\theta,[r])=2(c-b)\max a.
\]
These optimize over all rough competitors; they are distinct from a
Gauduchon variational volume with a prescribed probability.

For completeness the required small-curvature property has an explicit
construction. Put $\chi(t)=1-(1-t)^\alpha$, $\eta_L=\chi'(-L)$, and
\[
 \chi_L(t)=
 \begin{cases}
 \chi(t)-\eta_LL-\chi(-L),&t<-L,\\
 \eta_Lt,&-L\leq t\leq0.
 \end{cases}
\]
Values and derivatives agree at $-L$, and $\chi_L$ is convex increasing
with $0\leq\chi_L'\leq\eta_L\to0$. Hence the genuine function
$U_L=\chi_L(G_p)$ satisfies $U_L\sim U$ and
$\ddc U_L\geq-\eta_Lw_1$. The inequality extends canonically across
$D$. Thus $bG_p+U_L$ is an actual strict downstairs representative once
$\eta_L<c-b$.

\begin{lemma}\label{lem:pullback}
For a finite composition $\pi:Y'\to Y$ of point blowups,
$\pi^*U$ has zero point Lelong numbers everywhere, and its ordinary type
has arbitrarily small negative curvature on $Y'$.
\end{lemma}
\begin{proof}
The actual pullbacks $\pi^*U_L$ have bounded difference from $\pi^*U$,
including their poles. Fix $h'>0$ on $Y'$. Since the smooth semipositive
$\pi^*w_1$ is bounded by $C h'$, their Hessians are at least
$-C\eta_Lh'$.
To check all point Lelong numbers directly, write locally
$\pi^*G_p=\log|g|^2-\sigma$, where $g$ is a nonzero holomorphic
equation for the total transform of $D$. At a point of this divisor,
the germ has finite order $N$. Choose a direction $\xi$ with leading
homogeneous term $g_N(\xi)\ne0$; then
$g(t\xi)=t^N(g_N(\xi)+O(t))$. Along this direction $\pi^*U$ is
bounded below by $-C|\log|t||^\alpha-C$. Its spherical supremum has
the same lower bound. After a smooth local psh correction, the Lelong
growth quotient is therefore zero, since $\alpha<1$. Away from the
divisor the function is smooth. This argument concerns the specified
holomorphic-logarithm family and makes no assertion about pullbacks of
all zero-Lelong quasi-psh functions.
\end{proof}

\begin{theorem}\label{thm:blowup}
For every nonempty finite composition of point blowups $\pi:Y'\to Y$,
including centers in the pole set or an earlier exceptional divisor,
there is no strict $\pi^*\theta$-current with potential ordinarily
equivalent to $\pi^*r$. There is a strict witness
$v'\leq\pi^*r+C$. The two ordinary-type masses remain the displayed
positive values.
\end{theorem}
\begin{proof}
Apply Theorem~\ref{thm:criterion} with logarithmic divisor
$b\pi^*D$, residual $\pi^*U$ and background $\pi^*\theta$.
The resulting smooth residual background is $\pi^*\beta$.
The last exceptional curve $E\simeq\PP^1$ is contracted by the
composite $\pi$. If $\pi^*\beta+\ddc f>0$ for smooth $f$, restriction
to $E$ gives $\ddc(f|_E)>0$, because $\pi^*\beta|_E=0$.
Its integral is zero by curve Stokes, a contradiction. This proves the
mutual-type obstruction even when the original potential has a pole on
$E$: only the smooth $f$ has been restricted.

Construct a one-sided strict witness inductively. Suppose $v_j$ is strict
with lower Hermitian bound $\delta_jh_j$ on the current surface. For the
next point blowup $\pi_{j+1}$ with exceptional divisor $E_{j+1}$, choose
curvature $B_{E_{j+1}}$ of $\mathcal O(E_{j+1})$ negative on its vertical
tangent and a genuine $\ell_{E_{j+1}}$ with
$\ddc\ell_{E_{j+1}}=[E_{j+1}]-B_{E_{j+1}}$.
For sufficiently small $t_j>0$,
$\pi_{j+1}^*h_j-t_jB_{E_{j+1}}>0$.
Indeed in blowup charts $(x,y)\mapsto(x,xy)$ the pullback has a positive
horizontal coefficient, mixed coefficient $O(|x|)$ and Schur complement
at least $C_1|x|^2$. The added vertical curvature is at least $t_jC_2$;
its bounded mixed errors are $O(t_j|x|+t_j|x|^2+t_j^2)$.
Choose the exceptional neighborhood small, then $t_j$ small. On its
compact complement positivity persists under a bounded small perturbation.
Thus
\[
 v_{j+1}=\pi_{j+1}^*v_j+\delta_jt_j\ell_{E_{j+1}}
\]
has a strict current. Start with $v_0=bG_p+U_L$ as above. Genuine
quasi-psh pullbacks and their sums are locally integrable and canonical,
including over pole centers. The divisor logarithms are globally bounded
above, so induction gives $v'\leq\pi^*r+C$. The first part of the
theorem excludes mutual ordinary equivalence for this witness.

Finally, use the checked pure-pullback competitor bijection and actual
non-pluripolar pushforward identity of \cite{Ruled}. The bijection concerns
every competitor in the pulled-back ordinary type, and products carry no
mass on exceptional analytic sets. It preserves both infimum and supremum
of mass. The exceptional logarithms in the one-sided witness specify a
different ordinary type; their mass is not substituted for this optimization.
\end{proof}

\begin{remark}
The new proposed extension removes the outside-pole hypothesis in the
checked blowup obstruction of \cite{Ruled}, within the explicitly finite
point-blowup class. It supplies neither a factorization theorem for arbitrary
maps nor a strict-type criterion for every Hermitian ordinary singularity.
The small-curvature ordinary-type hypothesis implies zero point Lelong
numbers by the trace argument above; no additional zero-Lelong assumption
is imposed. General S11 remains open.
\end{remark}

\begin{thebibliography}{9}
\bibitem{Demailly92} J.-P.~Demailly,
\emph{Regularization of closed positive currents and Intersection Theory},
J.~Algebraic Geom. \textbf{1} (1992), 361--409,
Main Theorem~1.1 and Section~2.
\href{https://www-fourier.univ-grenoble-alpes.fr/~demailly/manuscripts/regularization.pdf}{Author-hosted primary text}.
\bibitem{DemaillyBook} J.-P.~Demailly,
\emph{Complex Analytic and Differential Geometry}, I.5.23--5.24, III.5.5--5.8 and III.6.9.
\href{https://www-fourier.univ-grenoble-alpes.fr/~demailly/manuscripts/agbook.pdf}{Author-hosted primary text}.
\bibitem{DemaillyBergman} J.-P.~Demailly,
\emph{On the cohomology of pseudoeffective line bundles},
arXiv:1401.5432v2, Theorem~1.2 and Section~2, especially Corollary~2.3.
Strong openness is due to Guan--Zhou; the supplied proof follows Hiep
and imports Ohsawa--Takegoshi and standard basis theory.
\href{https://arxiv.org/abs/1401.5432}{Primary text}.
\bibitem{Gauduchon} P.~Gauduchon,
\emph{Les m\'etriques standard d'une surface complexe compacte \`a premier
nombre de Betti pair}, Ast\'erisque \textbf{126} (1985), 129--135,
Theorem~2 and its proof, pp.~132--133.
\href{https://www.numdam.org/item/AST_1985__126__129_0/}{Published primary text}.
\bibitem{Ruled} \emph{Restricted masses and ordinary strict-current types
on a nonclosed ruled surface}, checked restricted construction, October~2026,
Sections~2 and~5. General S11 remains open.
\href{https://github.com/MingchenXia/non-kahler-pluripotential-theory/blob/2ec45dec6172ced6e2af18b2b5fe657c07aeee9e/research/papers/S11-nonclosed-ruled-restricted-masses-and-strict-types.tex}{Retained checked source}.
\end{thebibliography}
\end{document}
